\documentclass[12pt,a4paper]{article}

%% paquetes básicos
\usepackage[utf8]{inputenc}
\usepackage[spanish]{babel}
\usepackage{a4wide,color,verbatim,Sweave,url,xargs,amsmath,booktabs,longtable}
\usepackage{graphicx,float}
\usepackage{tikz,xcolor}
\usepackage{enumitem}

%% bibliotecas TikZ
\usetikzlibrary{automata,positioning,calc,arrows}

%% entornos para exams
\newenvironment{question}{}{}
\newenvironment{solution}{}{}
\newenvironment{answerlist}{\begin{enumerate}}{\end{enumerate}}

\begin{document}
\Sconcordance{concordance:ahorro_interpretacion_representacion_n3_v1.tex:ahorro_interpretacion_representacion_n3_v1.Rnw:1 %
20 1 1 240 53 1 2 2 6 1 2 2 6 1 2 2 6 1 2 2 6 1 1 80 10 0 1 4 13 1 1 2 %
4 0 1 2 8 1}



\begin{question}

Tres estudiantes implementaron diferentes estrategias de ahorro durante 6 meses. La siguiente tabla muestra el dinero ahorrado mensualmente por cada uno (en miles de pesos):

\par\vspace{0.5cm}

\begin{center}
\textbf{Ahorro mensual por estudiante (miles de pesos)}
\par\vspace{0.3cm}
\begin{tabular}{|c|c|c|c|}
\hline
\textbf{Mes} & \textbf{Mónica} & \textbf{Xavier} & \textbf{Ofelia} \\
\hline
1 & 40 & 30 & 50 \\
\hline
2 & 45 & 55 & 65 \\
\hline
3 & 70 & 75 & 80 \\
\hline
4 & 65 & 85 & 95 \\
\hline
5 & 85 & 90 & 110 \\
\hline
6 & 100 & 95 & 125 \\
\hline
\end{tabular}
\end{center}

\par\vspace{0.5cm}

Adicionalmente, se presenta la siguiente información calculada:

\begin{itemize}
\item Ahorro total de Mónica: 405 mil pesos
\item Ahorro total de Xavier: 430 mil pesos
\item Ahorro total de Ofelia: 525 mil pesos
\item Tasa de crecimiento promedio mensual de Mónica: 30\%
\item Tasa de crecimiento promedio mensual de Xavier: 43.3\%
\item Tasa de crecimiento promedio mensual de Ofelia: 30\%
\end{itemize}

\par\vspace{0.5cm}

Un director de banco necesita identificar cuál estudiante logró ahorrar la mayor cantidad total de dinero para otorgar un reconocimiento. Requiere una representación que muestre claramente los montos finales acumulados.

\par\vspace{0.5cm}

¿Cuál de las siguientes gráficas representa MEJOR los montos totales ahorrados por cada estudiante?

\par\vspace{0.5cm}

\textbf{A.}

\includegraphics{ahorro_interpretacion_representacion_n3_v1-002}

\textbf{porque muestra directamente los montos totales finales que cada estudiante logró acumular}

\par\vspace{0.5cm}

\textbf{B.}

\includegraphics{ahorro_interpretacion_representacion_n3_v1-003}

\textbf{porque solo compara dos puntos sin mostrar el total acumulado completo}

\par\vspace{0.5cm}

\textbf{C.}

\includegraphics{ahorro_interpretacion_representacion_n3_v1-004}

\textbf{porque permite visualizar las tendencias pero no los totales finales de forma clara}

\par\vspace{0.5cm}

\textbf{D.}

\includegraphics{ahorro_interpretacion_representacion_n3_v1-005}

\textbf{porque muestra porcentajes que no reflejan los montos reales acumulados}

\end{question}

\begin{solution}

Para identificar quién ahorró la mayor cantidad total, es necesario comparar directamente los montos finales acumulados por cada estudiante. \par\vspace{0.5cm}Análisis de cada gráfica presentada:\begin{itemize}\item \textbf{Gráfica A (Totales Acumulados) - CORRECTA}: muestra claramente los montos totales finales que cada estudiante logró acumular  porque:
  \begin{itemize}
  \item Presenta de forma clara y directa los montos totales finales
  \item Facilita la identificación inmediata del mayor ahorrador
  \item No requiere cálculos adicionales para la comparación
  \end{itemize}

\item \textbf{Gráfica B (Inicial vs Final)}: facilita la comparación directa y simple entre el punto de partida y el resultado final , pero  solo muestra dos puntos sin revelar información sobre eficiencia relativa o evolución temporal .

\item \textbf{Gráfica C (Evolución Mensual)}: permite visualizar las tendencias, fluctuaciones y patrones mes a mes , pero  requiere análisis adicional para comparaciones específicas ya que muestra mucha información temporal .

\item \textbf{Gráfica D (Tasas de Crecimiento)}: normaliza las diferencias iniciales y permite comparar directamente la eficiencia relativa , pero  muestra porcentajes que pueden no reflejar los montos reales o patrones temporales .\end{itemize}
\par\vspace{0.5cm}

Basándose en las tasas de crecimiento calculadas:
\begin{itemize}
\item Mónica: 30\% mensual
\item Xavier: 43.3\% mensual
\item Ofelia: 30\% mensual
\end{itemize}

\par\vspace{0.5cm}

Xavier tuvo la estrategia más eficiente en términos de crecimiento relativo.

\begin{answerlist}
  \item Verdadero
  \item Falso
  \item Falso
  \item Falso
\end{answerlist}
\end{solution}

%% META-INFORMATION
%% \extype{schoice}
%% \exsolution{1000}
%% \exname{Ahorro Interpretacion Representacion N3 V1}

\end{document}
